% !TeX root = ../../金融经济学_熊和平_习题.tex

\begin{problem}
\problemno{15.1} 什么是家庭金融？如何理解家庭金融从金融学中分离出来是必要的？
\end{problem}
\begin{solution}
家庭金融研究家庭在储蓄、借贷、住房、保险、退休和投资中的决策及其与金融市场的互动。家庭同时拥有劳动收入、人力资本、住房和非交易资产，并面临生命周期、家庭成员和制度约束，因此不能直接套用只有可交易金融资产的代表性投资者模型，独立研究有助于解释真实资产配置和福利差异。
\end{solution}
\begin{problem}
\problemno{15.2} 简述家庭金融理论与异质性投资者理论的关系，家庭作为经济人有哪些异质性？
\end{problem}
\begin{solution}
家庭金融把异质性投资者框架应用于家庭决策。家庭在财富、年龄、劳动收入风险、风险偏好、金融知识、健康状况、婚姻结构、税收和借贷约束上都可能不同，这些异质性造成储蓄率、住房和股票参与率差异。
\end{solution}
\begin{problem}
\problemno{15.3} 为什么说家庭资源禀赋和面临的经济环境更加复杂？
\end{problem}
\begin{solution}
家庭资源不仅有金融财富，还包括劳动收入、人力资本、住房和耐用品；支出还受子女、教育、医疗和养老需求影响。收入和支出风险具有长期性、不可交易性和相关性，且家庭受到税制、社会保险、信贷限额和信息成本约束，所以其预算集远比单一证券投资者复杂。
\end{solution}
\begin{problem}
\problemno{15.4} 什么是模糊性？简述风险、不确定性和模糊性的区别与联系。
\end{problem}
\begin{solution}
风险是结果集合和概率都已知；不确定性是结果可能已知但客观概率无法确定；模糊性强调对概率分布本身存在多种合理判断。三者都涉及未来状态，风险是最可量化的情形，模糊性可看作对概率的二阶不确定，通常会产生额外的模糊厌恶。
\end{solution}
\begin{problem}
\problemno{15.5} 家庭投资决策有哪些错误？如何理解这些错误？
\end{problem}
\begin{solution}
教材归纳的四类典型错误是：不参与资本市场从而放弃股票市场的福利改进；参与后分散不足、在相同风险下收益偏低；风险分担安排过于简单、承担不必要风险；住房按揭与再融资模式不科学（如高息负债与低息存款并存）。理解这些错误需结合有限理性、金融知识、社会规范与产品复杂性，而非仅归因于“不理性”。
\end{solution}
\begin{problem}
\problemno{15.6} 什么是基于生命周期的投资模型？
\end{problem}
\begin{solution}
生命周期模型把家庭的劳动收入、消费、储蓄和投资放在多个年龄阶段联合决定。年轻时人力资本高、金融财富少，通常积累资产；中年收入和财富较高，进行风险分散；退休后依靠金融资产和养老金消费。最优股票比例取决于剩余寿命、劳动收入风险、风险厌恶和退休制度，而非固定不变。
\end{solution}
\begin{problem}
\problemno{15.7} 某家庭为了购买住房在银行办理了 \num{50} 万元按揭贷款，贷款利率为 $6\%$，同时他们在其他银行拥有 \num{50} 万元定期存款，存款利率为 $3\%$。当问及为何购买住房使用贷款时，答曰：他熟悉的朋友在买房时都使用了贷款，何况银行给他们的贷款利率比一般的贷款利率优惠。请分析该家庭该项金融决策的错误。
\end{problem}
\begin{solution}
该家庭同时借款和持有同额存款，相当于以 $6\%$ 的利率借入资金再以 $3\%$ 的利率投资，忽略税收和流动性差异时每年确定损失为 $500000(6\%-3\%)=15000$ 元。决策依据“朋友都这样做”是从众和熟悉性偏差，优惠贷款利率低于一般贷款并不能说明低于其存款机会成本。正确比较应考虑提前还贷收益、存款流动性、税收、违约风险和住房的非金融收益；若无特殊约束，先用存款偿还高利率贷款更有效率。
\end{solution}
